\sectionkicker{Source-backed technical notes}
\subsection*{巨大化を成立させる裏側――ZeROとGShard: Technical Notes}
本欄は記事本文で圧縮した一次資料上の情報を、比較・再検証しやすい形で整理したものである。時系列、確認済みの主張、留意点、一次資料URLを掲載する。
\smallskip
\noindent{\small\textit{Technical Notesでは、一次資料から確認した公開・提供・研究上の事実を記載する。数値・能力評価や提供元・著者による比較表現は、独立再現済みの結果ではなく帰属claimとして扱う。}}\par
\medskip
\subsection*{Theme at a glance}
\addcontentsline{toc}{subsection}{Theme at a glance}
\begin{center}
\footnotesize
\begin{tabularx}{\linewidth}{@{}p{0.20\linewidth}p{0.10\linewidth}p{0.14\linewidth}X@{}}
\toprule
資料 & 位置づけ & 種別 & 時系列 \\
\midrule
DeepSpeed / ZeRO & 主要資料 & Framework & 2020-02-13 \\
GShard & 主要資料 & 論文 & 2020-06-30 \\
\bottomrule
\end{tabularx}
\end{center}
\normalsize

\begin{technicalnote}{DeepSpeed / ZeRO}{主要資料}
\begingroup
% reader-facing Technical Notes break policy
\widowpenalty=10000
\clubpenalty=10000
\displaywidowpenalty=10000
\begin{tabularx}{\linewidth}{@{}>{\bfseries}p{0.22\linewidth}X@{}}
組織 & Microsoft \\
種別 & Framework \\
時系列 & 2020-02-13 \\
\end{tabularx}
\smallskip
\begin{minipage}{\linewidth}
% reader-facing Technical Notes event-only fact/source tail group
{\bfseries 一次資料から整理したtechnical points}
\begin{itemize}[leftmargin=1.5em,itemsep=0.35em]
\item \textbf{一次情報で確認できる事実}: Microsoftの選定済み一次資料では、ZeROはdata-parallel process間でoptimizer states・gradients・parametersを複製せずpartitionするmemory optimizationとして説明され、3段階がそれぞれoptimizer state、gradient、parameterのpartitioningに対応する。 2020-02-13の公開時点でDeepSpeedに実装・releaseされたのはZeRO stage 1のoptimizer-state partitioningで、stage 2とstage 3は将来対応として記載されている。 DeepSpeedはPyTorch-compatibleなopen-source distributed-training libraryとして公開され、100B-parameter model supportやthroughput改善値はMicrosoft側の測定・主張として扱う。
\end{itemize}
\begin{samepage}
{\bfseries 一次資料}
\begingroup\sloppy
\Urlmuskip=0mu plus 2mu\relax
\begin{itemize}[leftmargin=1.5em,itemsep=0.25em]
\item {\scriptsize\url{https://www.microsoft.com/en-us/research/blog/zero-deepspeed-new-system-optimizations-enable-training-models-with-over-100-billion-parameters}}
\end{itemize}
\endgroup
\end{samepage}
\end{minipage}
\endgroup
\end{technicalnote}
\medskip

\begin{technicalnote}{GShard}{主要資料}
\begingroup
% reader-facing Technical Notes break policy
\widowpenalty=10000
\clubpenalty=10000
\displaywidowpenalty=10000
\begin{tabularx}{\linewidth}{@{}>{\bfseries}p{0.22\linewidth}X@{}}
組織 & Google / authors \\
種別 & 論文 \\
時系列 & 2020-06-30 \\
\end{tabularx}
\smallskip
\begin{minipage}{\linewidth}
% reader-facing Technical Notes event-only fact/source tail group
{\bfseries 一次資料から整理したtechnical points}
\begin{itemize}[leftmargin=1.5em,itemsep=0.35em]
\item \textbf{一次情報で確認できる事実}: Google / authorsの選定済み一次資料では、GShardはlightweight annotation APIsとXLA compiler extensionでparallel computation patternを表現し、既存model codeへの変更を抑えながらautomatic shardingを行うsystemとして提示された。 著者らはSparsely-Gated Mixture-of-Expertsを使うmultilingual Transformerを600B parameter超へscaleし、2048 TPU v3 acceleratorsで4日trainingしたと報告する。 model scale・training time・translation qualityはこの論文の著者実験に属し、一般的なdistributed-training性能保証とはみなさない。
\end{itemize}
\begin{samepage}
{\bfseries 一次資料}
\begingroup\sloppy
\Urlmuskip=0mu plus 2mu\relax
\begin{itemize}[leftmargin=1.5em,itemsep=0.25em]
\item {\scriptsize\url{http://arxiv.org/abs/2006.16668v1}}
\end{itemize}
\endgroup
\end{samepage}
\end{minipage}
\endgroup
\end{technicalnote}
\medskip
